\documentclass{article}
\usepackage{amsmath,amssymb}
\begin{document}
\section*{Two-stage intermodal network design with disruption scenarios}
This is the complete selected model for the supplied instances. The original
root-level extraction and the source paper are preserved. Arc notation
$e=(i,j,m)\in E$ combines node pair and mode; $N,K,S$ denote nodes, orders and
scenarios, with probabilities $\pi_s$. Order $k$ has origin $o_k$, destination
$d_k$, quantity $r_k$ and due time $T_k$. All times use the input's common time
unit. Arc parameters are capacity $q_e$, vehicle capacity $v_e$, distance
$D_e$, nominal time $t_e$, fixed cost $f_e$, variable cost $c_e$, consolidation
cost $c^c_e$, emissions $e_e,e^c_e$, social cost per unit-distance $c^s_e$
and road vehicle externality per distance $c^r_e$ (zero for other modes).
Global parameters are budget $A$, carbon price $C$, lost-demand cost $u$,
late-time cost $\hat u$, transfer cost $c^t$, emission $e^t$ and time $\tau^t$.
Scenario inputs are capacity reduction $g_{es}$ and time increase $\delta_{es}$.
The supplied action sets contain one preparation and one recovery option per
arc: costs $h_e=20000$, $a_e=5000$, discounted recovery cost
$\tilde a_e=4000$ after preparation, capacity gains $0.2q_e,0.5q_e$,
recovered travel time $\hat t_{es}=t_e(1+0.3\delta_{es})$, and implementation
time $o_e=2$. These are instance parameters for the corresponding actions.

\subsection*{Variables and objective}
First-stage fleet $Y_e\in\mathbb Z_{\ge0}$ and preparation
$\beta_e\in\{0,1\}$ are common to every scenario.
Recourse variables are $X_{kes},U_{ks},A_{ks},DL_{ks}\ge0$,
$Z_{kes},\gamma_{es}\in\{0,1\}$, and
$Y'_{es},Y^k_{es},n_{kjs},\operatorname{con}_{es}\in\mathbb Z_{\ge0}$.
$A_{ks}$ is the arrival-time upper certificate at $d_k$. $n$ is transferred quantity,
measured in transferred units.
\begin{align*}
\min\quad&
\sum_e(f_e+c^r_eD_e)Y_e+\sum_eh_e\beta_e\\
&+\sum_s\pi_s\left[
\sum_{k,e}\{c_e+D_e(Ce_e+c^s_e)\}X_{kes}
+\sum_{k,j\ne d_k}(c^t+Ce^t)n_{kjs}
+\sum_e(c^c_e+Ce^c_e)\operatorname{con}_{es}
+\sum_e a_e\gamma_{es}
+\sum_k(uU_{ks}+\hat u DL_{ks})\right].
\end{align*}
\subsection*{Complete constraints}
For every applicable order, node, arc and scenario:
\begin{align}
\sum_{e:\,\mathrm{tail}(e)=j}X_{kes}
-\sum_{e:\,\mathrm{head}(e)=j}X_{kes}
&=\begin{cases}r_k-U_{ks}&j=o_k\\-r_k+U_{ks}&j=d_k\\0&\text{otherwise},\end{cases}\\
0\le U_{ks}&\le r_k,\\
\sum_kX_{kes}&\le(1-g_{es})q_e+0.2q_e\beta_e+0.5q_e\gamma_{es},\\
\sum_kX_{kes}&\le v_eY'_{es},\qquad X_{kes}\le v_eY^k_{es},\\
\operatorname{con}_{es}&=\sum_kY^k_{es}-Y'_{es},\qquad Y'_{es}\le Y_e,\\
Z_{kes}&\le X_{kes}\le r_kZ_{kes},\\
\sum_{\substack{e:\,\mathrm{head}(e)=j\\\mathrm{mode}(e)=m}}X_{kes}
-\sum_{\substack{e:\,\mathrm{tail}(e)=j\\\mathrm{mode}(e)=m}}X_{kes}
&\le n_{kjs}\qquad(j\ne d_k,\text{ each mode }m),\\
\sum_e Z_{kes}t_e(1+\delta_{es})
+\sum_e Z_{kes}\gamma_{es}\{\hat t_{es}-t_e(1+\delta_{es})\}
+\sum_e o_e\gamma_{es}+\sum_{j\ne d_k}\tau^t n_{kjs}
&\le A_{ks},\\
DL_{ks}&\ge A_{ks}-T_k,\qquad DL_{ks}\ge0,\\
\sum_e h_e\beta_e+\sum_ea_e\gamma_{es}
+\sum_e(\tilde a_e-a_e)\beta_e\gamma_{es}&\le A,\\
\gamma_{e0}&=0.
\end{align}
Binary products can be represented by their standard exact binary
linearization. One action of each type is available per arc; hence its
binary domain implements the corresponding mutual-exclusion constraint.

\subsection*{Explicit corrections and algorithm boundary}
The source paper's domain (18) makes $Y$ a nonnegative integer; the old
extracted binary domain was wrong. The source and extraction both print
$A-T_kDL\le DL$ although the text defines $DL$ as delay time. That product
is represented here by the time-consistent delay bound $DL\ge A-T_k$ above.
The delay definition uses consistent time units. The code and all rerun objectives use this
corrected late-time definition. The source's unusual global recovery
implementation-time sum is retained literally.

The solver constructs an integer-flow feasible candidate on the complete
network. Integer flows, at most two preparations and at most four recoveries,
shortest-path choices, minimal vehicle counts and minimal modal transfer
quantities are search choices only. They guide candidate construction.
Every original node, arc, order and scenario remains in the certificate.
The returned objective is a feasible upper bound.
\end{document}
